\documentclass[11pt,a4paper]{article}
\usepackage[T1]{fontenc}
\usepackage[utf8]{inputenc}
\usepackage{lmodern}
\usepackage{microtype}
\usepackage[a4paper,margin=2.25cm]{geometry}
\usepackage{longtable,booktabs,array,calc}
\usepackage{fancyvrb}
\usepackage{xcolor}
\usepackage{hyperref}
\usepackage{xurl}
\usepackage{fancyhdr}
\usepackage{enumitem}
\usepackage{parskip}
\usepackage{titlesec}
\definecolor{ddtadark}{RGB}{30,38,48}
\definecolor{ddtablue}{RGB}{38,82,126}
\hypersetup{colorlinks=true,linkcolor=ddtablue,urlcolor=ddtablue,citecolor=ddtablue}
\pagestyle{fancy}
\fancyhf{}
\lhead{Documentation-Driven Threat Analysis}
\rhead{BA3 work plan R15}
\cfoot{\thepage}
\setlength{\headheight}{14pt}
\setlist[itemize]{leftmargin=1.5em,itemsep=0.2em,topsep=0.25em}
\setlist[enumerate]{leftmargin=1.7em,itemsep=0.2em,topsep=0.25em}
\titleformat{\section}{\Large\bfseries\color{ddtadark}}{}{0em}{}
\titleformat{\subsection}{\large\bfseries\color{ddtadark}}{}{0em}{}
\titleformat{\subsubsection}{\normalsize\bfseries\color{ddtadark}}{}{0em}{}
\providecommand{\tightlist}{\setlength{\itemsep}{0pt}\setlength{\parskip}{0pt}}
\setlength{\emergencystretch}{4em}
\hbadness=10000
\hfuzz=20pt
\sloppy
\begin{document}
\section{DDTA research work plan after documentation-layer
closure}\label{ddta-research-work-plan-after-documentation-layer-closure}

\textbf{WORK PLAN - REVISION 15}

\textbf{Status:} active BA3 execution plan after BA3-T3 derivation
reproducibility/change-impact pressure testing.

\textbf{Supersedes:} Revision 14 only for forward execution state;
R1-R14 remain historical research records.

\subsection{1. Fixed prior state}\label{fixed-prior-state}

\begin{itemize}
\tightlist
\item
  Chapters 2-4: CLOSED / FINAL for current thesis scope.
\item
  Documentation layer: CLOSED.
\item
  BA0-R systems-modeling prior-art research: CLOSED.
\item
  BA0 responsibility/non-goals: CLOSED.
\item
  BA1 minimal BAE identity ontology: CLOSED.
\item
  BA2 relation/action vocabulary: CLOSED BY BA2-T4.
\item
  BA3-T1 provenance/origin lower bound: COMPLETED / PROVISIONAL PASS.
\item
  BA3-T2 identity/lifecycle pressure test: COMPLETED / PROVISIONAL PASS
  WITH IDENTITY-LIFECYCLE REFINEMENT.
\item
  \texttt{BAReferent} and \texttt{BAProposition}: ACCEPTED first-class
  semantic identity families.
\item
  ThreatForge is a case study/tooling instrument, not DDTA semantic
  authority.
\end{itemize}

\subsection{2. BA3 objective}\label{ba3-objective}

BA3 defines the smallest provenance, derivation, identity/lifecycle and
change-revalidation mechanics that connect governed documentation to
Base Analysis without creating a second source of project truth.

\subsection{3. BA3-T1 - provenance/origin lower
bound}\label{ba3-t1---provenanceorigin-lower-bound}

\textbf{Status: COMPLETED / PROVISIONAL PASS WITH LOWER-BOUND
CANDIDATE.}

T1 requires provenance independently on BAReferent/BAProposition,
baseline-scoped source locators, many-to-many lineage, explicit origin
state and derivation basis/rationale for derived meaning.

\subsection{4. BA3-T2 - cross-baseline identity, staleness and
lifecycle}\label{ba3-t2---cross-baseline-identity-staleness-and-lifecycle}

\textbf{Status: COMPLETED / PROVISIONAL PASS WITH IDENTITY-LIFECYCLE
REFINEMENT.}

T2 closes the provisional separation among origin, review/freshness and
identity continuity. It retains
\texttt{RETAIN\ \textbar{}\ REPLACE\ \textbar{}\ RETIRE}, with accepted
replacement/retirement preserving historical
\texttt{SUPERSEDED}/\texttt{RETIRED} semantics.

\subsection{5. BA3-T3 - derivation-rule reproducibility and
change-impact
lineage}\label{ba3-t3---derivation-rule-reproducibility-and-change-impact-lineage}

\textbf{Status: COMPLETED / PROVISIONAL PASS WITH DERIVATION-IMPACT
REFINEMENT.}

T3 pressure-tested facial M4 plus provider-state normalization and the
analysis-to-governance feedback chain.

\subsubsection{Main dispositions}\label{main-dispositions}

\begin{verbatim}
plain untyped derivationBasis                 REJECTED
role-bound derivationBasisBinding             REQUIRED
immutable inspectable rule revision           REQUIRED
universal executable derivation language      REJECTED
semantic replay reproducibility               REQUIRED
explicit revalidationContext                  REQUIRED
revalidationContext == BA2 dependOn           REJECTED
M4 silent BA repair from sibling Decision     REJECTED
localized staleness propagation               REQUIRED
whole-BA invalidation on any source change    REJECTED
provider mapping hidden in tool               REJECTED
analysis result/candidate as BA source         REJECTED
new BAE family                                 NOT FORCED
BA1 / BA2 / T1 / T2 reopen                    NOT TRIGGERED
\end{verbatim}

\subsubsection{Active refined provenance
shape}\label{active-refined-provenance-shape}

\begin{verbatim}
BAOriginAttachment
|- targetElement             1
|- governedBaselineKey       1
|- originState               1
|- sourceLink                0..*
|- derivationBasisBinding    0..*
|    |- inputRoleKey         1
|    `- basisRef             1
|- derivationRuleRef         0..1
`- revalidationContext       0..*
\end{verbatim}

A derivation rule reference resolves to an immutable method-neutral
registry revision with role, applicability and output semantic
contracts. The registry is extensible; no universal exhaustive rule list
is required.

\subsection{6. Why BA3 remains open}\label{why-ba3-remains-open}

T3 resolves the remaining pressure targets but intentionally does not
self-close BA3. The accumulated T1/T2/T3 candidate now needs one
integrated adversarial closure review to check:

\begin{itemize}
\tightlist
\item
  redundancy among source, derivation, context and continuity metadata;
\item
  contradictions among origin/review/freshness/lifecycle dimensions;
\item
  cross-corpus sufficiency under M1-M4 and provider/responsibility
  changes;
\item
  authority preservation through the analysis/corrective feedback loop;
\item
  whether any new responsibility actually forces BA1/BA2 reopen.
\end{itemize}

\subsection{7. BA3-T4 - provenance, identity/lifecycle and
derivation/change-impact closure
review}\label{ba3-t4---provenance-identitylifecycle-and-derivationchange-impact-closure-review}

\textbf{NEXT / NOT EXECUTED BY THIS PLAN REVISION.}

BA3-T4 must perform only the integrated closure review.

It must pressure-test at least:

\begin{enumerate}
\def\labelenumi{\arabic{enumi}.}
\tightlist
\item
  whether any T1/T2/T3 metadata responsibility can be removed or merged
  without losing source drill-down, reproducibility, uncertainty
  localization or targeted revalidation;
\item
  whether \texttt{originState}, \texttt{reviewState}, \texttt{freshness}
  and continuity disposition remain orthogonal enough to justify
  separate dimensions;
\item
  whether role-bound derivation basis plus immutable rule descriptor is
  sufficient across the current derived examples without a universal
  derivation language;
\item
  whether \texttt{revalidationContext} is necessary and sufficiently
  narrow, especially for M4, without leaking into BA2 project semantics;
\item
  whether the localized staleness rules over-invalidate or
  under-invalidate the facial M1-M4 and order/provider controls;
\item
  whether the feedback chain preserves governed documentation as the
  only project authority after accepted/rejected corrections;
\item
  whether the combined BA3 candidate forces a new first-class BAE
  family, new BA2 operator or another material reopen;
\item
  whether BA3 can close for current thesis scope with explicit future
  reopen criteria.
\end{enumerate}

\subsubsection{BA3-T4 guardrails}\label{ba3-t4-guardrails}

Do not:

\begin{itemize}
\tightlist
\item
  start BA4 projections;
\item
  define formal STRIDE/STRIDE-AI overlay schemas;
\item
  define AnalysisRecord/Common Finding;
\item
  design ThreatForge classes/tables;
\item
  require one graph/storage/serialization;
\item
  turn revalidation metadata into general systems-model dependencies;
\item
  demand exhaustive derivation-rule enumeration absent concrete
  evidence.
\end{itemize}

\subsubsection{BA3-T4 exit condition}\label{ba3-t4-exit-condition}

If no material counterexample survives, close BA3 for current thesis
scope and freeze the smallest accepted provenance/derivation/lifecycle
contract with explicit reopen criteria. Otherwise reopen only the
smallest failed BA3 responsibility and leave BA3 open.

\subsection{8. BA4 - projections}\label{ba4---projections}

Only after BA3 closes, define derived human and method projections from
the same accepted BA semantics. No view becomes project authority.

\subsection{9. BA5 - lexical vocabulary and optional
assistance}\label{ba5---lexical-vocabulary-and-optional-assistance}

Only after BA3/BA4 boundaries are stable, define display labels,
authoring synonyms and optional assistance. NLP/LLM proposals remain
candidate-producing support.

\subsection{10. BA6 - complete Base Analysis regression and
closure}\label{ba6---complete-base-analysis-regression-and-closure}

Regress the complete BA0-BA5 design against the closed corpora and at
least one structurally different holdout. BA6 may reopen only the
smallest earlier responsibility on a material counterexample.

\subsection{11. Analysis envelope remains
downstream}\label{analysis-envelope-remains-downstream}

Only after Base Analysis closure:

\begin{itemize}
\tightlist
\item
  A1 - AnalysisRecord / execution envelope;
\item
  A2 - common reviewed Finding boundary;
\item
  A3 - change/provenance integration;
\item
  formal methodology overlays and STRIDE/STRIDE-AI evaluations.
\end{itemize}

\subsection{12. Next authorized
microstep}\label{next-authorized-microstep}

Only after the BA3-T3 package is reviewed, committed, pushed and
remotely verified, execute:

\begin{quote}
\textbf{BA3-T4 - provenance, identity/lifecycle and
derivation/change-impact closure review.}
\end{quote}

Do not start BA4 or downstream analysis schemas before BA3 is explicitly
closed.

\end{document}
